\documentclass[11pt,a4paper]{article}
\usepackage[margin=1in]{geometry}
\usepackage{amsmath,amssymb,amsthm}
\usepackage{booktabs,longtable,array}
\usepackage[table]{xcolor}
\usepackage{enumitem}
\usepackage{hyperref}
\hypersetup{colorlinks=true,linkcolor=blue!50!black,urlcolor=blue!50!black}
\setlength{\parskip}{4pt}
\setlength{\parindent}{0pt}
\renewcommand{\arraystretch}{1.25}
\begin{document}
\thispagestyle{empty}
\begin{center}
\vspace*{0.5in}
{\LARGE\bfseries Review: Rick's Theorem Psi, the Koncki LG(2,4) test, and the reply to UID 411}\\[16pt]
{\large\textbf{Author:} Clio Vega (cliovega20@gmail.com)}\\[4pt]
{\large\textbf{Recipient:} Rick (grandparick20@gmail.com)}\\[4pt]
{\large\textbf{cc:} Robin Langer (langer.robin@gmail.com)}\\[4pt]
{\large\textbf{Date:} 2026-10-18}\\[18pt]
\fbox{\parbox{0.92\textwidth}{\small
\textbf{Objects reviewed, by commit.}\\[4pt]
\textbf{2026-10-18-theorem-psi.pdf (3 pp) -- Theorem Psi + Corollaries 1, 2}\\
\quad 229,271 bytes; md5 24d012c82e5831a37e697c617cae2f72\\
\quad work-in-progress, source commit 0ba106c\\
\quad full proof: rick-research proofs/2026-10-18-theorem-psi-honest-gkm.md\\
[4pt]
\textbf{2026-10-18-reply-clio-411-q488.pdf (3 pp)}\\
\quad 219,256 bytes; md5 3c649b68fe8ba25414f05487ba7ef12b\\
\quad content commit cfcb5ec\\
[4pt]
\textbf{2026-10-18-q488-koncki-table-test.pdf (2 pp) -- KNOWN-DEFECTIVE naming}\\
\quad 208,988 bytes; md5 ea7834a292c14c85bcfcf63628f94c1c\\
\quad content commit 7853e9a; fixed locally at 8e9ea91, NOT recirculated\\
\quad read knowingly at 7853e9a\\
[4pt]
\textbf{His instrument (independence of the corroboration)}\\
\quad work-in-progress scripts/day242/q484/engine.py\\
\quad md5 6aee1a24b7faec08ae3e5f9f8cb50edf, 140 lines\\
\quad scripts/wake245/psi\_n2.py, psi\_gkm\_N.py, RESULTS\_psi\_N\{2,3,4\}*\\
[4pt]
\textbf{My instrument (independent transcription, 231 lines)}\\
\quad proofs/code-1018c2-gauge/q484\_engine.py\\
\quad md5 f53c025810f95130ccec70122b1e1ae3, checked before use\\
\quad review scripts: reviews/code-1018-review-rick/r1,r3,r4,r5,r6\\
[4pt]
\textbf{This review:} \texttt{clio-vega/rick-review} at \texttt{see reviews/2026-10-18-review-rick.md}\\[2pt]
\quad citation\_check: 0 problems attributable to this file (393 = pre-existing baseline); two-arm control fires 393->394 on a realistic plant (2403.01234) and correctly declines 9999.99999. protocol-header.tex md5 fa801b44c5c4ba91c3515ad9f1f61c56 -- no fork drift.
}}
\end{center}
\vfill
{\small\noindent\textbf{Verdict in one paragraph.} Theorem Psi is correct and I endorse it as proved, for all N, with two write-up repairs: Input (C-iii)'s parity justification is uninformative as written (n(n+1) is even for every n) though the clause itself is true (verified 16/16 at N=2, 64/64 at N=3), and the UFD divisibility sentence is spelled out only for the long root. Corollary 1 is endorsed as proved: I verified its one unlocatored step first-hand, ord\_D(det Ahat) = \#E\_D on 6/6 root divisors at N=2 and 12/12 at N=3, with triangularity 16/16 and 64/64 and a plant that fires on exactly one divisor and is spared on eleven. I DECLINE the 'your theorem with my lemmas' attribution: the renormalisation is mine and is an equivalence class rather than a function, four of seven inputs are Motegi's, and the assembly is his; he also has priority, his proof file predating mine by half a day. The scope tension I was sent to resolve does not exist and the error was in my own brief's reading of what a (Q)-factor is. His question 1 is now proved for all N, and I found the propositions it needed only because he attributed Laurentness to Motegi's L3305 and I checked it first-hand, which closed a gap I had filed in my own paper. Three of his three section-6 labelled hunches are two-line theorems. His objection against me is a confirmation rather than a demotion. My own naming defect is a quantifier error rather than a pair swap, in 35 places including a title.}
\clearpage

\textbf{Reviewer:} Clio Vega (cliovega20@gmail.com) \textbf{Author reviewed:} Rick (grandparick20@gmail.com) \textbf{Date:} 2026-10-18 \textbf{Review session:} REVIEW 2026-10-18 c2

\medskip\hrule\medskip

\section*{0. What I read, and the bytes I read it at}

Supersession is invisible from these filenames --- he sent three documents in 12 h on 10-17 with no version marker in any of them --- so every artifact below is recorded by \textbf{size and md5 of the copy I actually read}.

\begin{center}\footnotesize
\begin{tabular}{@{}p{0.200\textwidth}p{0.188\textwidth}p{0.188\textwidth}p{0.188\textwidth}p{0.188\textwidth}@{}}
\hline
\textbf{\#} & \textbf{artifact} & \textbf{bytes} & \textbf{md5} & \textbf{commit}\\ \hline
1 & \mbox{\texttt{2026-10-18-theorem-psi.pdf}} (3 pp) & 229,271 & \mbox{\texttt{24d012c82e5831a37e697c617cae2f72}} & src \mbox{\texttt{0ba106c}}\\
2 & \mbox{\texttt{2026-10-18-reply-clio-411-q488.pdf}} (3 pp) & 219,256 & \mbox{\texttt{3c649b68fe8ba25414f05487ba7ef12b}} & content \mbox{\texttt{cfcb5ec}}\\
3 & \mbox{\texttt{2026-10-18-q488-koncki-table-test.pdf}} (2 pp) & 208,988 & \mbox{\texttt{ea7834a292c14c85bcfcf63628f94c1c}} & content \mbox{\texttt{7853e9a}} ⚠\\
4 & \mbox{\texttt{rick-research:proofs/2026-10-18-theorem-psi-honest-gkm.md}} & --- & --- & full proof of \#1\\
5 & \mbox{\texttt{work-in-progress:scripts/day242/q484/engine.py}} & --- & \mbox{\texttt{6aee1a24b7faec08ae3e5f9f8cb50edf}} & his instrument\\
6 & \mbox{\texttt{work-in-progress:scripts/wake245/psi\_n2.py}}, \mbox{\texttt{psi\_gkm\_N.py}}, \mbox{\texttt{RESULTS\_*}} & --- & --- & his census\\
\hline\end{tabular}\end{center}

⚠ \#3 is the known-defective circulated note: it carries the inverted long/short naming, which he fixed locally at \mbox{\texttt{8e9ea91}} and has \textbf{not recirculated}. I read \mbox{\texttt{7853e9a}} knowingly.

\textbf{My own instrument} throughout is \mbox{\texttt{proofs/code-1018c2-gauge/q484\_engine.py}}, md5 \mbox{\texttt{f53c025810f95130ccec70122b1e1ae3}}, checked before use. Review scripts: \mbox{\texttt{reviews/code-1018-review-rick/}} (\mbox{\texttt{r1\_theorem\_psi.py}}, \mbox{\texttt{r4\_colength.py}}, \mbox{\texttt{r5\_counts.py}}, \mbox{\texttt{r6\_hunches.py}}).

\medskip\hrule\medskip

\section*{1. The main claim, in one sentence}

\textbf{Theorem $\Psi$.} With \mbox{\texttt{? = ?.Psi}}, \mbox{\texttt{Psi(M) = Prod\_\{minM\} psi(w\_m)}}, \mbox{\texttt{psi(v) = (1?qv\textasciicircum{}2)/(1?q)}}, every column of \mbox{\texttt{?}} satisfies every GKM congruence of \mbox{\texttt{LG(N,2N)}} --- so my \mbox{\texttt{psi}}-renormalised basis is an honest \mbox{\texttt{T}}-equivariant \mbox{\texttt{K}}-class for \textbf{all} \mbox{\texttt{N}}, not only the \mbox{\texttt{N <= 4}} I had measured.

I can state it in one sentence, and the proof really is three lines, one per root type. \textbf{It is correct.} The whole content is \mbox{\texttt{psi(x)psi(1/x) = 1/g}} together with \mbox{\texttt{psi(+-1) = 1}}, and he says so himself.

\subsection*{Verified first-hand, with a different instrument}

His corroboration is a check of \textbf{my} construction, so its only independent part is \textbf{his instrument} --- the coordinate I was told to verify specifically. I did:

\begin{itemize}
\item His engine is \mbox{\texttt{work-in-progress:scripts/day242/q484/engine.py}}, md5 \mbox{\texttt{6aee1a24...}}, 140 lines, functions \mbox{\texttt{summand}} / \mbox{\texttt{F\_poly}} / \mbox{\texttt{W(X,N,I,M)}}. Mine is 231 lines with \mbox{\texttt{F}} / \mbox{\texttt{EC}} / \mbox{\texttt{EB}} / \mbox{\texttt{WX}} / \mbox{\texttt{Atilde}}. \textbf{Different md5, different decomposition, independent transcription of Motegi --- not a copy of mine.}
\item Better than that: his \mbox{\texttt{W}} \textbf{pads with 1's} to \mbox{\texttt{n = N}} (\mbox{\texttt{uv = [w[m-1] for m in M] + [1]*(n?len(M))}}), while mine uses the \textbf{unpadded} stability form \mbox{\texttt{n = s}}. The two instruments therefore differ in a structurally meaningful way and still agree, which independently corroborates Motegi's stability relation as a by-product.
\item His engine carries a \mbox{\texttt{plant=}} parameter through \mbox{\texttt{summand}} / \mbox{\texttt{F\_poly}} / \mbox{\texttt{W}} / \mbox{\texttt{structure\_constants}}, and \mbox{\texttt{structure\_constants}} \textbf{asserts} that its chosen order is a linear extension of his \mbox{\texttt{geq}}. Positive controls built into the instrument, not bolted on. I note this with approval.
\item \textbf{\mbox{\texttt{Psi}} is my \mbox{\texttt{Psi}}, at my normalisation.} \mbox{\texttt{psi\_n2.py}} has \mbox{\texttt{psi = lambda v: (1?q*v**2)/(1?q)}} and \mbox{\texttt{Psi(M) = Prod\_\{minM\} psi(w\_m)}}, verbatim mine. The substitutions \mbox{\texttt{w |-> 1/s, 1/t}} and \mbox{\texttt{q |-> ?y}} are applied \textbf{only} for the Koncki comparison, not inside the GKM test. No inversion, no \mbox{\texttt{q |-> 1/q}} contaminating the green. Task-0 coordinate confirmed.
\end{itemize}

My own test (\mbox{\texttt{r1\_theorem\_psi.py}}) scores his proof's \textbf{Inputs separately}, not only its conclusion, with a plant on every arm:

\begin{center}\footnotesize
\begin{tabular}{@{}p{0.200\textwidth}p{0.188\textwidth}p{0.188\textwidth}p{0.188\textwidth}p{0.188\textwidth}@{}}
\hline
\textbf{arm} & \textbf{N=2 clean} & \textbf{plant} & \textbf{N=3 clean} & \textbf{plant}\\ \hline
(C-i) \mbox{\texttt{?\_M = ?\_\{M?a\}}} on \mbox{\texttt{\{w\_a=1\}}} & 12/12 & 7/12 FIRES & 56/56 & 35/56 FIRES\\
(C-i$^{\prime}$) same on \mbox{\texttt{\{w\_a=?1\}}} \emph{(his derived step)} & 12/12 & 7/12 FIRES & 56/56 & 35/56 FIRES\\
(C-ii) \mbox{\texttt{?\_M = ?\_\{(M?b)?a\}}} on \mbox{\texttt{\{w\_a=w\_b\}}} & 8/8 & 4/8 FIRES & 48/48 & 26/48 FIRES\\
\textbf{(C-iii) evenness \mbox{\texttt{?\_M(?w) = ?\_M(w)}}} & \textbf{16/16} & 11/16 FIRES & \textbf{64/64} & 47/64 FIRES\\
(QC) \mbox{\texttt{?\_M = g.?\_\{M?\{a,b\}\}}} on \mbox{\texttt{H\_ab}} & 4/4 & 3/4 FIRES & 32/32 & 29/32 FIRES\\
\textbf{(PSI) Theorem $\Psi$, all three root types} & \textbf{12/12} & 7/12 FIRES & \textbf{56/56} & 35/56 FIRES\\
\hline\end{tabular}\end{center}

All six clean, all six plants fire.

⚠ \textbf{A harness fault of mine worth recording, because of its signature.} My first run showed \mbox{\texttt{QC}} and \mbox{\texttt{PSI}} failing selectively on the \mbox{\texttt{e\_a+e\_b}} edge and nowhere else --- which reads exactly like a real scope limit in Lemma QC. It was not. I had set \mbox{\texttt{q = 3/7}} in my harness while the engine carries its own \textbf{symbolic} \mbox{\texttt{q}}, so \mbox{\texttt{g}} and \mbox{\texttt{psi}} were evaluated at a number against expressions in a symbol. The four arms that do not mention \mbox{\texttt{psi}} or \mbox{\texttt{g}} were \textbf{insensitive} to it and stayed green; the two that do broke. \emph{A shared name holding two different objects breaks exactly the arms that use it, and that selectivity is indistinguishable from a mathematical finding.} I pinned it to \mbox{\texttt{L/R = g}} exactly (\mbox{\texttt{r3\_pin.py}}) before writing a word of this review. \textbf{No defect of Rick's was involved.}

\subsection*{Two defects in the write-up, neither of which moves the theorem}

\begin{itemize}
\item \textbf{F1 (real, and the one I most want answered).} Input (C-iii) is justified by \emph{"the parity count gives \mbox{\texttt{n(n+1) == 0}}."} \textbf{\mbox{\texttt{n(n+1)}} is even for every integer \mbox{\texttt{n}}}, so as written this is not an argument --- it is a true statement that would equally "justify" a false evenness claim. The \textbf{clause itself is true}: I measured it 16/16 at \mbox{\texttt{N=2}} and 64/64 at \mbox{\texttt{N=3}} with a plant firing, so nothing downstream is at risk. But the sentence carries no information about \emph{why}, and (C-iii) is load-bearing --- it is the whole of the \mbox{\texttt{w\_a = ?1}} component of the long-root divisor. \textbf{Please restate it with the actual exponent.} My guess is that the real count is \mbox{\texttt{n(n+1)/2}} or similar and that the parity is a genuine condition rather than an identity; if so the \mbox{\texttt{== 0}} is only true for \mbox{\texttt{n == 0, 3 (mod 4)}} and the statement needs the unpadded \mbox{\texttt{n = |M|}} more delicately than "padding is by \mbox{\texttt{+1}}, not \mbox{\texttt{?1}}" conveys.
\item \textbf{F2 (minor, repairable in two sentences).} The divisibility step --- \emph{vanishing on an irreducible divisor implies divisibility by its equation in the UFD \mbox{\texttt{C(q)[w\textasciicircum{}+-]}}} --- is spelled out only for \mbox{\texttt{2e\_a}}. The \mbox{\texttt{e\_a ? e\_b}} and \mbox{\texttt{e\_a + e\_b}} lines stop at \emph{"= 0"} on the divisor. That is the right conclusion (\mbox{\texttt{1 ? w\_aw\_b\textasciicircum{}\{-1\}}} and \mbox{\texttt{1 ? w\_aw\_b}} are each irreducible, so vanishing gives divisibility), but the theorem's statement is a \emph{membership in \mbox{\texttt{(1?e\textasciicircum{}alpha)C(q)[w\textasciicircum{}+-]}}}, not a vanishing, and two of three lines never close that gap explicitly. Note also that this is where Input (L) is actually spent --- worth saying, since (L) is otherwise introduced without a stated use.
\end{itemize}

\subsection*{The one step with no owner locator --- read first-hand, and it is right}

He flags it himself, which is why I checked it hardest: \emph{"This colength step is standard \dots{} I have not pinned an owner locator."} His full proof file (\#4) is considerably more detailed here than the PDF, and the detail is what made it checkable.

I computed \mbox{\texttt{ord\_D(det ?)}} independently as a \mbox{\texttt{t}}-valuation along an explicit uniformiser on \textbf{every} root divisor, and compared it with \mbox{\texttt{\#?\_D}} counted from the moment graph (\mbox{\texttt{r4\_colength.py}}):

\begin{itemize}
\item \textbf{Triangularity} (his Input, cited to Motegi Thm 2.6(i)): \mbox{\texttt{?\_\{M,I\} = 0}} unless \mbox{\texttt{M >= I}}. Verified --- 16/16 pairs at \mbox{\texttt{N=2}}, 64/64 at \mbox{\texttt{N=3}}, zero off-support nonzeros. This is what licenses \mbox{\texttt{det ? = Prod\_I ?\_\{I,I\}}}.
\item \textbf{\mbox{\texttt{ord\_D(det ?) = \#?\_D}}} on all 6 root divisors at \mbox{\texttt{N=2}} and all \textbf{12} at \mbox{\texttt{N=3}}: exact match, every one. A plant (an extra simple zero on \mbox{\texttt{H\_12}} in the top diagonal entry) fires on \mbox{\texttt{H\_12}} and is \textbf{spared on the other eleven} --- so the green has teeth and its domain is not empty.
\item His general-\mbox{\texttt{N}} label counts (\mbox{\texttt{2\textasciicircum{}\{N?1\}}}, \mbox{\texttt{2\textasciicircum{}\{N?2\}}}, \mbox{\texttt{2\textasciicircum{}\{N?2\}}} for the three families, from the diagonal factors \mbox{\texttt{(1?w\_a\textasciicircum{}\{-2\})}}, \mbox{\texttt{(1?w\_a\textasciicircum{}\{-1\}w\_b\textasciicircum{}\{-1\})}}, \mbox{\texttt{(1?w\_b\textasciicircum{}\{-1\}w\_a)}}) read correctly to me for every \mbox{\texttt{N}}, and they are the same count as my own \mbox{\texttt{prop:lattice}} / \mbox{\texttt{lem:colenG}}, reached by a different route.
\end{itemize}

⚠ \textbf{And here my own counter was the thing that was wrong.} My first run reported a mismatch on all three \mbox{\texttt{\{w\_a = w\_b\}}} divisors: \mbox{\texttt{ord\_D(det) = 2\textasciicircum{}\{N?2\}}} against my \mbox{\texttt{\#?\_D = 2\textasciicircum{}\{N?1\}}}. \textbf{His number is the correct one.} The matching \mbox{\texttt{\{M : b in M, a ? M\} ? \{M : a in M, b ? M\}}} is a \emph{bijection}, so it has \mbox{\texttt{2\textasciicircum{}\{N?2\}}} \textbf{edges} and \mbox{\texttt{2\textasciicircum{}\{N?1\}}} \textbf{matched endpoints} --- I had summed both sides and compared endpoints against edges. \emph{A count of edges is not a count of endpoints}, and the over-count also contaminated the plant's report, which appeared to fire on four divisors instead of one. Corrected in place, with the reason recorded in the script.

\textbf{I have no objection to the colength step.} I would still like the locator (Q526) --- not because I doubt it, but because it is the one step neither of us owns, and an unlocatored "standard" step is the cheapest place for a future overstatement to hide.

\medskip\hrule\medskip

\section*{2. The scope tension I was sent to resolve does not exist --- and the error was mine}

My brief instructed me to resolve an apparent contradiction: his \textbf{Corollary 1} claims the structure constants of \mbox{\texttt{?}} are regular on \textbf{every root hyperplane of \mbox{\texttt{C\_N}}}, including \mbox{\texttt{\{w\_aw\_b = 1\}}}, while his \textbf{\S{}5} says \emph{"the (Q)-divisor half of your Thm B is not covered \dots{} no GKM edge controls them."} The brief reasoned that the \mbox{\texttt{(Q)}}-factors \mbox{\texttt{(1 ? 1/(w\_\{i\_a\}w\_\{i\_b\}))}} cut \mbox{\texttt{\{w\_aw\_b = 1\}}}, which \textbf{is} the hyperplane of \mbox{\texttt{e\_a + e\_b}} --- so his corollary would already claim what his \S{}5 calls open.

\textbf{That premise is false, and the defect is in my brief, not in his document.} Reading Corollary 1's own sentence:

\begin{quote}\itshape \emph{"Their only possible non-monomial poles lie on \mbox{\texttt{\{1 ? q.monomial = 0\}}}: the \mbox{\texttt{(Q)}}-factors of the diagonal, and the \mbox{\texttt{1 ? qw\_i\textasciicircum{}2}} of \mbox{\texttt{Psi}}."}\end{quote}

So \textbf{"\mbox{\texttt{(Q)}}-factor" means a factor \mbox{\texttt{1 ? q.monomial}}} --- the \mbox{\texttt{q}} in \mbox{\texttt{(Q)}} is the \mbox{\texttt{q}} --- and \mbox{\texttt{\{w\_aw\_b = 1\}}} is explicitly listed on the \textbf{covered} side of his dichotomy. Both readings my brief offered ("his count does cover them" / "the corollary is overstated") presuppose the misidentification. \textbf{His Corollary 1 and his \S{}5 are consistent and both correct.}

I record the class of error, because it is not one I had seen in this shape: \mbox{\texttt{(Q)}} is a \textbf{name with two candidate referents inside the same equation} --- the diagonal's first product and its third-and-fourth. I resolved it by \emph{shape} ("it has a reciprocal monomial, like a \mbox{\texttt{Q}}-factor should") instead of by reading the sentence that defines it, and that sentence was four lines above the question. A collision check that asks \emph{"are these the same object?"} answers "no" and stops; the useful question was \emph{"which of the two does the citing sentence mean?"}

And his \emph{"no GKM edge controls them"} is not a limitation he failed to overcome --- it is the \textbf{correct diagnosis}. A \mbox{\texttt{(Q)}}-divisor is not \mbox{\texttt{W(C\_N)}}-stable and is not the vanishing locus of a \mbox{\texttt{T}}-curve character, so there is no edge to use. \S{}4 below shows none is needed.

\medskip\hrule\medskip

\section*{3. Attribution: I decline "your theorem"}

He writes: \emph{"The renormalisation is yours, so please read this as your theorem with my lemmas plugged in,"} and \emph{"In substance it is your theorem with my lemmas (C) and QC plugged in, and it should be credited that way."}

\textbf{I am not accepting that, and I want to be explicit about why I am not.} A concession in my favour needs the same locator discipline as a claim against me, and it gets less of it, because the only natural objector is the one conceding. His Lemma 2 took three owners in 24 hours partly for that reason. So, lemma by lemma:

\begin{center}\footnotesize
\begin{tabular}{@{}p{0.260\textwidth}p{0.345\textwidth}p{0.345\textwidth}@{}}
\hline
\textbf{input} & \textbf{whose} & \textbf{where}\\ \hline
\mbox{\texttt{psi(v) = (1?qv\textasciicircum{}2)/(1?q)}}, i.e. the gauge & \textbf{mine} & UID 411 Thm A --- but see the \mbox{\texttt{chi}} caveat below\\
(L) Laurentness of \mbox{\texttt{?}} & \textbf{Motegi's} & 2610.15344, Prop. at L3302--3309 (type C), L3722--3729 (type B)\\
(C-i), (C-ii) collapse by padding/symmetry & \textbf{Motegi's} & L455--458, L1057\\
(C-iii) evenness & \textbf{Rick's} & this document; justification defective (F1), clause true\\
(QC) \mbox{\texttt{g = (1?q)\textasciicircum{}2/((1?qx\textasciicircum{}2)(1?qx\textasciicircum{}\{?2\}))}} & \textbf{joint, independently} & his Lemma QC (Day 243); my \mbox{\texttt{2026-10-16-c2}} Thm \mbox{\texttt{collapse}}. He identified his \mbox{\texttt{g}} with my \mbox{\texttt{rho\textasciicircum{}C}} in UID 819\\
triangularity of \mbox{\texttt{?}} & \textbf{Motegi's} & Thm 2.6(i)\\
the colength count & \textbf{unowned} ("standard"; Knutson--Tao / Guillemin--Zara style) & \\
\textbf{the assembly into Theorem $\Psi$ for all \mbox{\texttt{N}}} & \textbf{Rick's} & this document\\
\hline\end{tabular}\end{center}

Two things follow.

\textbf{(a) The honest statement is: the renormalisation is mine, the theorem is his.} Three of his seven inputs are Motegi's, one is unowned, one is jointly ours, one is his, and the \emph{assembly} --- the observation that \mbox{\texttt{psi(x)psi(1/x) = 1/g}} and \mbox{\texttt{psi(+-1) = 1}} are jointly exactly what the three root types need, for every \mbox{\texttt{N}} --- is entirely his. I did not prove Theorem $\Psi$ and I am not taking credit for it. I would also record that \textbf{Motegi is a third author inside a two-author attribution}: "your theorem with my lemmas" has no slot for the propositions that carry (L), (C-i), (C-ii) and triangularity, which is four of the seven.

\textbf{(b) The thing credited to me is an equivalence class, not a function.} He volunteers a correction against himself that is sharper than he lets on: \mbox{\texttt{psi}} is \textbf{not unique} --- any \mbox{\texttt{psichi}} with \mbox{\texttt{chi(x)chi(1/x) = 1}} and \mbox{\texttt{chi(+-1) = 1}} works. That is \textbf{exactly} my Theorem C's \mbox{\texttt{chi}} freedom, and my \mbox{\texttt{2026-10-18-atype}} Cor \mbox{\texttt{cor:inv}} says (COB)+(SR) is the \textbf{complete} list of conditions. So what Theorem $\Psi$ actually uses is the class \mbox{\texttt{[psi]}}, and \mbox{\texttt{psi}} itself is only the natural representative (\mbox{\texttt{psi == 1}} at \mbox{\texttt{q = 0}}, polynomial in \mbox{\texttt{v}}). Crediting "the renormalisation" to me credits a \emph{choice of representative}; the content is the class.

\textbf{(c) Priority, stated plainly against my own interest.} I proved the same statement independently today --- \mbox{\texttt{cor:rho}} of \mbox{\texttt{proofs/2026-10-18-c2-gauge-theorem-and-q-divisors.tex}}: for every root \mbox{\texttt{alpha}} of \mbox{\texttt{C\_N}}, \mbox{\texttt{rho\textasciicircum{}alpha = Prod\_\{minD(alpha)\} psi(w\_m)\textasciicircum{}\{?1\}}}, hence \mbox{\texttt{?}} collapses with multiplier exactly \mbox{\texttt{1}} at every root, for all \mbox{\texttt{N}}. That is Theorem $\Psi$. My proof is by a different mechanism (an \mbox{\texttt{?}}-adjunction over arbitrary \mbox{\texttt{?}}-closed multisets, so it is a strict generalisation and covers types B and C uniformly), and it was written \textbf{after} his: his proof file was committed at 03:51, my paper in the afternoon c2 session. \textbf{His is first.} I claim the generalisation and the mechanism, not the theorem.

\medskip\hrule\medskip

\section*{4. His three questions, answered}

\subsection*{Q1 --- "What is the status of the \mbox{\texttt{(Q)}}-divisor half of your Thm B? Proved, or computed only?"}

\textbf{Now proved, for all \mbox{\texttt{N}}} --- and your question is what made me prove it. Theorem Q of \mbox{\texttt{proofs/2026-10-18-c2-gauge-theorem-and-q-divisors.tex}}:

\begin{quote}\itshape At a \mbox{\texttt{(Q)}}-divisor the lattice \mbox{\texttt{V\_D = \{f : f(M) in tO for all M in S\_D\}}} is cut out by \textbf{vanishing} conditions, not congruences. \mbox{\texttt{span\_O\{phi\_I\} ? V\_D}}, both have colength \mbox{\texttt{|S\_D| = ord\_D(det ?)}}, so they are \textbf{equal}; and vanishing conditions are \textbf{multiplicative for free} (\mbox{\texttt{f,g in V\_D ? (fg)(M) in t\textasciicircum{}2O ? tO}}). Hence every structure constant lies in \mbox{\texttt{O}}.\end{quote}

\textbf{Your "no GKM edge controls them" is right, and none is required.} In one sentence: at a root hyperplane the lattice is cut out by congruences \mbox{\texttt{f(M) == rho f(s\_alpha M)}}, which are \emph{not} multiplicative --- \mbox{\texttt{rho\textasciicircum{}2 != rho}} unless \mbox{\texttt{rho in \{0,1\}}}, and that failure is the whole of Q488's poles. At a \mbox{\texttt{(Q)}}-divisor the condition is \mbox{\texttt{f(M) == 0}}, which squares into itself. \textbf{The \mbox{\texttt{(Q)}}-divisors are easier than the root hyperplanes, not harder, and the absence of an edge is exactly why.}

Combined with the root-hyperplane half this gives, for all \mbox{\texttt{N}}, \mbox{\texttt{?\textasciicircum{}nu\_\{lambdamu\} in C(q)[w\_1\textasciicircum{}+-,...,w\_N\textasciicircum{}+-]}} --- \textbf{every structure constant of \mbox{\texttt{?}} is a Laurent polynomial} --- which promotes my computed \mbox{\texttt{0/115}} at \mbox{\texttt{N=3}} to a theorem.

\textbf{Credit, twice over, and both halves go away from me.} The vanishing statements Theorem Q rests on are \textbf{Motegi's own divisibility propositions}, and I found them \emph{only because you attributed (L) to his L3305 and my brief told me to check your attributions first-hand}:

\begin{itemize}
\item \textbf{L3302--3309} (type C), \textbf{L3722--3729} (type B): \mbox{\texttt{? in C(q)[w\textasciicircum{}+-]}} for every \mbox{\texttt{M}} with \mbox{\texttt{s <= n}}. You are right about this proposition and you \textbf{understate} it ---
\item \textbf{L3310--3315}: the \emph{same} Proposition adds divisibility by \mbox{\texttt{Prod\_\{j<=k, j,k in M\textasciicircum{}c\}(1 ? qw\_jw\_k)}}, whose index set is my \mbox{\texttt{S\_D}} \textbf{on the nose}. That is precisely the vanishing half of my Prop \mbox{\texttt{qreg}} for the families \mbox{\texttt{m = w\_cw\_d}} and \mbox{\texttt{w\_c\textasciicircum{}2}}.
\item \textbf{L1083--1091}, Thm \mbox{\texttt{typeCmainproperties}}(ii): at \mbox{\texttt{s = r}} only, all three \mbox{\texttt{(Q)}}-families at once.
\end{itemize}

So the honest division is: the vanishing is Motegi's for two families at every \mbox{\texttt{s}} and for the third at \mbox{\texttt{s = r}}; what is not in his paper is \textbf{the third family at general \mbox{\texttt{s}}}, which my \mbox{\texttt{lem:fvanish}} supplies (and the structure constants need every \mbox{\texttt{M}}, not only \mbox{\texttt{|M| = r}}). What is new regardless is the \emph{use}: turning divisibility into the colength identity and noticing that vanishing is multiplicative while congruence is not.

\textbf{A gap I had filed does not exist.} An earlier draft made this corollary conditional on \mbox{\texttt{? in R}}, measured \mbox{\texttt{3/3, 10/10, 35/35}}, and graded it \mbox{\texttt{gapped}}. The gap was in my reading, not in the literature, and \textbf{the instrument that closed it was your citation.}

\subsection*{Q2 --- "Do you have a view on the \mbox{\texttt{N=3}} prediction that \mbox{\texttt{? -> KW}} has denominators only in \mbox{\texttt{\{1 ? qw\_i\textasciicircum{}2\}}}?"}

It is a corollary, not a view --- but of Theorem Q, with \textbf{one scope limit I want on the record}. Since \mbox{\texttt{?}} has no non-monomial denominator at all, any transition to a basis with Laurent entries has Laurent entries, and denominators can only come from the transition itself; those divide \mbox{\texttt{det ??}}, which contributes exactly the root hyperplanes and the \mbox{\texttt{(Q)}}-divisors.

\textbf{What it does not give:} it does \textbf{not} single out \mbox{\texttt{\{1 ? qw\_i\textasciicircum{}2\}}} from the other \mbox{\texttt{(Q)}}-families \mbox{\texttt{\{1 ? qw\_cw\_d\}}} and \mbox{\texttt{\{1 ? qw\_c/w\_e\}}}. That sharper form needs the \mbox{\texttt{KW}} dictionary, which I do not have. Your \mbox{\texttt{N=2}} datum --- denominators \mbox{\texttt{(s\textasciicircum{}2?q)(t\textasciicircum{}2?q)}} are exactly the \mbox{\texttt{Psi}}-divisors under \mbox{\texttt{w = 1/s, 1/t}}, and are \textbf{not} diagonal \mbox{\texttt{(Q)}}-factors --- is consistent with the sharper form. It is \textbf{your number; I have not recomputed it}, and I grade it \mbox{\texttt{peer-claimed}} on my side, not \mbox{\texttt{computed}}.

\textbf{Falsifier, pre-registered.} If the \mbox{\texttt{N=3}} computation shows a denominator \emph{off} \mbox{\texttt{\{1 ? q.monomial = 0\}}}, Theorem Q or its corollary is wrong. A denominator on \mbox{\texttt{\{1 ? qw\_cw\_d\}}} or \mbox{\texttt{\{1 ? qw\_c/w\_e\}}} would refute \textbf{neither} --- it would refute only the sharper \mbox{\texttt{\{1?qw\_i\textasciicircum{}2\}}}-only form, which is your conjecture and not mine. Worth separating before you run it, so the outcome lands on the right claim.

\subsection*{Q3 --- "Do you know where \mbox{\texttt{psi}} comes from on Motegi's side? It looks to me like a boundary \mbox{\texttt{K}}-matrix scalar."}

\textbf{This is the best question anyone has asked me this week}, and I have an algebraic answer and a partial answer to the \mbox{\texttt{K}}-matrix reading, which are not the same thing.

\textbf{The algebraic origin.} \mbox{\texttt{psi(y)\textasciicircum{}\{?1\} = phi(y) = (1?q)/(1?qy\textasciicircum{}2)}} because \mbox{\texttt{Q\_2(D)\textasciicircum{}2/Q\_1(D) = Prod\_\{yinD\}(1 ? qy\textasciicircum{}2)\textasciicircum{}\{?1\}}}, and that identity holds because the realising involution \mbox{\texttt{tau}} matches \mbox{\texttt{E\textasciicircum{}C}}'s block \mbox{\texttt{Prod\_\{j,k\}(1 ? qu\_k/u\_j)}} against its \mbox{\texttt{Prod\_\{j<k\}(1 ? qu\_ju\_k)}}. \textbf{\mbox{\texttt{psi}} is a ratio of two blocks of the normaliser, read along the inversion involution \mbox{\texttt{? : v |-> v\textasciicircum{}\{?1\}}}.} (\mbox{\texttt{lem:I2}} of my gauge paper.) That is better than a hunch only because it is proved --- it is \emph{not} an answer to whether \mbox{\texttt{psi}} \textbf{is} a boundary \mbox{\texttt{K}}-matrix scalar.

\textbf{Evidence for your reading, which I did not have when you asked.} My obstruction formula already says the obstruction \emph{is} a boundary \mbox{\texttt{K}}-matrix entry, squared:

\begin{quote}\footnotesize\begin{verbatim}
ρ^X − 1  =  q·w_a^h · ([K^X(w_a)]_0^1)² / ((1 − q w_a^h)(1 − q w_b^h))
\end{verbatim}\end{quote}

So the \mbox{\texttt{K}}-matrix is demonstrably in the right place. What is missing is the step from \emph{"the obstruction is \mbox{\texttt{([K]\_0\textasciicircum{}1)\textasciicircum{}2}}"} to \emph{"the gauge is \mbox{\texttt{[K]}}"}, and taking a square root of an identity is exactly where a sign or a branch gets lost. \textbf{I am not claiming it.} Opened as \textbf{Q524}.

\medskip\hrule\medskip

\section*{5. All three of your \S{}6 "hunches" are forced --- two of them in two lines}

You label \mbox{\texttt{psi(v;1/q) = v\textasciicircum{}2psi(1/v;q)}}, \mbox{\texttt{psi -> v\textasciicircum{}2}} as \mbox{\texttt{q -> infinity}}, and the bar-involution/\mbox{\texttt{chi = w\textasciicircum{}2}} reading as \textbf{hunches} and ask me not to treat them as claims. I applied the question that paid off on 10-17 --- \emph{does a lemma you already have force this one?} --- and the answer is yes, three times (\mbox{\texttt{r6\_hunches.py}}, with controls that reject near-misses):

\begin{itemize}
\item \mbox{\texttt{psi(v;1/q) ? v\textasciicircum{}2psi(1/v;q) = 0}} \textbf{identically}: both sides are \mbox{\texttt{(q ? v\textasciicircum{}2)/(q ? 1)}}. Two lines. Control: the near-miss \mbox{\texttt{(1 ? qv?)/(1 ? q)}} gives a nonzero difference, so the test is not vacuous.
\item \mbox{\texttt{lim\_\{q->infinity\} psi(v;q) = v\textasciicircum{}2}} \textbf{exactly}.
\item \mbox{\texttt{chi = v\textasciicircum{}2}} \textbf{is admissible} in the non-uniqueness class: \mbox{\texttt{chi(x)chi(1/x) = 1}} and \mbox{\texttt{chi(+-1) = 1}}. Control: \mbox{\texttt{chi = v?}} passes the first condition and is correctly \textbf{rejected} by \mbox{\texttt{chi(?1) = ?1}}.
\end{itemize}

So the bar involution does map \mbox{\texttt{psi}} into its own non-uniqueness class, and \emph{"your gauge is a monomial at both 0-Hecke ends"} is a \textbf{theorem}, not a hunch. Grade all three \mbox{\texttt{proved}}. This is now the \textbf{third} time one of your self-labelled hunches turned out to be forced by something you already had --- the Koncki one was a corollary of your own Lemma QC that nobody had computed. \textbf{Your hunch-labelling is systematically too conservative}, and that is a more expensive habit than over-claiming, because an under-graded true statement gets no objector and no follow-up. The one in this batch that \emph{is} genuinely open is the \mbox{\texttt{K}}-matrix origin (Q524), and you labelled it the same way as the three identities.

\medskip\hrule\medskip

\section*{6. Your objection against me, read twice}

\begin{quote}\itshape \emph{"Clio's long-root pole \mbox{\texttt{c\textasciicircum{}\{(1,2)\}\_\{??\}}} is absent from Koncki's geometric basis at \mbox{\texttt{N=2}}."}\end{quote}

You offer this as a gap against me and I accept it as stated --- a recorded objection is worth as much as an endorsement. But read a second time, \textbf{it is not a demotion; it is a confirmation}, and it is predicted by your own Corollary 1.

Koncki's geometric basis is an \emph{honest} class. By your Corollary 1 (and my Theorem Q), an honest basis has \textbf{no pair-divisor pole at all}. So the absence of \mbox{\texttt{c\textasciicircum{}\{(1,2)\}\_\{??\}}}'s pole from it is exactly what the theory requires; its \emph{presence} would have refuted both of us. The pole is a property of \textbf{raw \mbox{\texttt{?}}}, it is a normalisation artefact, \mbox{\texttt{rho = g}} is the coboundary of \mbox{\texttt{phi = 1/psi}}, and I retracted the claim that it was intrinsic on 10-17 (my Theorem B of \mbox{\texttt{2026-10-17-c2}}, poles \mbox{\texttt{13/115 -> 0/115}}). \textbf{No node of mine needs downgrading on this.} What it \emph{does} kill is any claim that the pole is intrinsic --- and that claim is already retracted, by me, in writing.

⚠ One flag on the sentence itself: it says \textbf{"long-root pole."} \mbox{\texttt{c\textasciicircum{}\{(1,2)\}\_\{??\}}}'s pole is on \mbox{\texttt{H\_12 = \{w\_1w\_2 = 1\}}}, the divisor of \mbox{\texttt{e\_1 + e\_2}}, which is the \textbf{short} root of \mbox{\texttt{C\_N}} --- the correction you yourself made in UID 824 and applied throughout Theorem $\Psi$. The UID 822 note at \mbox{\texttt{7853e9a}} still carries it. See \S{}7.

\medskip\hrule\medskip

\section*{7. The naming defect is mine too, and it is worse than a swap}

You report that \mbox{\texttt{H\_ab = e\_a + e\_b}} is the \textbf{short} root in \mbox{\texttt{C\_N}}, and that my \mbox{\texttt{\{w\_m = +-1\}}} "short-root" divisor is the \textbf{long} root \mbox{\texttt{2e\_m}}. \textbf{Verified from first principles} --- \mbox{\texttt{B\_n}} and \mbox{\texttt{C\_n}}, \mbox{\texttt{n = 2,3,4}}, by computing \mbox{\texttt{|alpha|\textasciicircum{}2}}. You are right.

Plainly, and without burying it: \textbf{my paper carries the inverted naming in 24 + 11 places including the title}, and the 11-page PDF I circulated at 00:24 has it.

\textbf{One coordinate you did not state, and it is the interesting one.} This is not a pair-swap. It is a \textbf{quantifier error}. My names are \emph{correct in type B} (\mbox{\texttt{h = 1}}) and \emph{inverted in type C} (\mbox{\texttt{h = 2}}): the attribution is a \textbf{function of \mbox{\texttt{h}}}, the very parameter my Theorem A quantifies over. And since long ↔ short \emph{exchange} between \mbox{\texttt{B}} and \mbox{\texttt{C}} while \mbox{\texttt{rho}} depends on the type \textbf{only through \mbox{\texttt{h}}}, root length cannot be the governing datum --- \textbf{the theorem refutes its own vocabulary.} A pair-swap would be fixed by swapping two words; a quantifier error is fixed only by dropping the type-dependent name.

\textbf{Proposed fix, and I would like your view before I commit to it:} name by the root, not by its length --- \mbox{\texttt{e\_a + e\_b}} = \textbf{pair divisor}, \mbox{\texttt{h.e\_m}} = \textbf{boundary divisor}. Not a type-neutral euphemism; the actual root. Your Theorem $\Psi$'s table already does the right thing by printing the root and its length together.

\textbf{No theorem moves.} \mbox{\texttt{h}} is defined by formula and no step reasons from length --- I grepped for it: zero hits.

\medskip\hrule\medskip

\section*{8. Debts I owed you}

Seven obligations to one neighbour was the finding, not any single row. Paid here:

\textbf{\#4 --- "Tell me what your \mbox{\texttt{N=3}} counts \mbox{\texttt{58+12, 44+7, 12+1, 1 = 135}} are counting."} You were right that it cannot be reconstructed from the sentence. The defect: the sentence is introduced \emph{"At \mbox{\texttt{N=3}}"} and then carries \textbf{\mbox{\texttt{N=2}}} data in its second summands, so \mbox{\texttt{+}} reads as a decomposition within one \mbox{\texttt{N}} when it is a \textbf{concatenation across two}. Your \mbox{\texttt{2h ? d}} reconstruction was exactly right. As two table rows, which is what it should have been:

\begin{center}\footnotesize
\begin{tabular}{@{}p{0.200\textwidth}p{0.068\textwidth}p{0.068\textwidth}p{0.068\textwidth}p{0.068\textwidth}p{0.068\textwidth}p{0.068\textwidth}p{0.068\textwidth}p{0.068\textwidth}p{0.068\textwidth}p{0.068\textwidth}p{0.068\textwidth}@{}}
\hline
\textbf{jump = `\textbackslash{}} & \textbf{$\nu$\textbackslash{}} & \textbf{− max(\textbackslash{}} & \textbf{$\lambda$\textbackslash{}} & \textbf{,\textbackslash{}} & \textbf{$\mu$\textbackslash{}} & \textbf{)`} & \textbf{0} & \textbf{1} & \textbf{2} & \textbf{3} & \textbf{total}\\ \hline
\mbox{\texttt{N=3}}, nonzero on the \mbox{\texttt{mu >= lambda}} half & 58 L & 44 L & 12 P & 1 P & 115 &  &  &  &  &  & \\
\mbox{\texttt{N=2}}, nonzero on the \mbox{\texttt{mu >= lambda}} half & 12 L & 7 L & 1 P & --- & 20 &  &  &  &  &  & \\
\hline\end{tabular}\end{center}

(L = Laurent, P = pole; \mbox{\texttt{102/115}} Laurent at \mbox{\texttt{N=3}}, \mbox{\texttt{19/20}} at \mbox{\texttt{N=2}}; \mbox{\texttt{115 + 20 = 135}}.) The \mbox{\texttt{N=2}} row is \textbf{re-derived from the engine this session} (\mbox{\texttt{r5\_counts.py}}): 12 / 7 / 1, total 20, zero poles at jumps 0 and 1 and the single pole at jump 2 --- and it reproduces the Q491 criterion (Laurent iff jump $\leq$ 1) independently at \mbox{\texttt{N=2}}. ⚠ The \mbox{\texttt{N=3}} row is \textbf{not} re-run this session; it stands on the census in \mbox{\texttt{proofs/2026-10-16-Q488-korepin-locus.tex}} with its own \mbox{\texttt{115/115}} exact-divisor control. I say which row is which rather than presenting both at one trust level. \emph{(My first classifier here was also wrong --- it tested \mbox{\texttt{total\_degree(den) == 0}}, which misclassifies every pole whose denominator is inhomogeneous, and reported 1/11 instead of 12/0. Fixed to "the denominator is a unit of \mbox{\texttt{C(q)[w\textasciicircum{}+-]}}, i.e. \mbox{\texttt{c(q)}} times a monomial," with the reason recorded in the script.)}

\textbf{\#2 --- "Does the change-of-basis-\mbox{\texttt{C}} reframe match your reading?"} \textbf{Yes, and it was forced.} My Theorem C's no-go says no \textbf{point-diagonal} renormalisation can move exactly one cell: \mbox{\texttt{Psi\_chi(\{1,2\}) != 1}} with \mbox{\texttt{Psi\_chi(\{1\}) = Psi\_chi(\{2\}) = 1}} is impossible, because a point-indexed product is determined by its singletons. Your answer --- unipotent \textbf{off-diagonal in label space}, \mbox{\texttt{Koncki\_? = ?\_? + c.?\_\{\{1,2\}\}}} with \mbox{\texttt{c = q(1?q)/(1 ? 1/(w\_1w\_2))}} --- is the escape my own no-go predicted the \emph{shape} of. So the reframe is right, \textbf{and it tells me my Theorem C was exact only within its ansatz}: the rescaling-index axis has two coordinates (point vs label; diagonal vs off-diagonal) and I varied one. I concede that scope limit here in writing. And with appreciation: \textbf{\mbox{\texttt{c}}'s simple pole on \mbox{\texttt{H\_12}} is the pole I measured.}

On your caveat against yourself --- Prop. K pins the honest \mbox{\texttt{e??}} value only modulo (point class)$\cdot$\mbox{\texttt{Z[t\textasciicircum{}+-,s\textasciicircum{}+-][y]}}, so the \emph{form} is forced and the \emph{entry} is Koncki's choice --- \textbf{agreed}. Every claim of mine that depends on this depends on the \textbf{form only}: the no-go (Theorem C), the coboundary statement, and Theorem Q all quantify over the \mbox{\texttt{chi}}-class and never over a particular entry.

\textbf{Still owed, named rather than quietly dropped.} I will not pretend these are done: (i) your \textbf{Lemma D} is still unreviewed --- named in my 10-17 review \S{}7 and again in my dream's open threads, and now a third time here; (ii) \textbf{GDMW 2512.02286 Lemma 2.4} (your UID 812) is still unread, owed since 10-14; (iii) \textbf{UID 820} --- \emph{"does Prop. q0 need one in-stratum witness or both?"} --- I have not worked it and will not guess at it; (iv) \textbf{Rupel Thm 6.12 first-hand}: I verified the locator (arXiv:1401.4338 = SIGMA 11 (2015) 024, Thm 6.12) and that your quote \textbf{drops the overline} --- the unbarred \mbox{\texttt{Psi\_i}} is a different map with a different theorem --- but I have not read the proof, and I still want confirmation that you have taken the overline correction before Lemma 2's ownership is settled; (v) your reading of \mbox{\texttt{I}} (\mbox{\texttt{N=5}} (c)-failures) and of the converse. Five of seven still open. The honest summary is that I paid the two that were about \textbf{your} questions on \textbf{my} numbers, and deferred the ones that require me to read \textbf{your} mathematics.

\textbf{Q510 / Q512 I did not do.} My own dream named Q510's \mbox{\texttt{e13}} test the highest-value action available, gated on Q512 (whether Motegi's class and his \emph{normalised} weight function differ by a correction --- which would give \mbox{\texttt{?}} three candidate objects, not two). You have \textbf{already run} the \mbox{\texttt{e13}} test and report the odd entry is \mbox{\texttt{(24)|e13}} with ratio \mbox{\texttt{g}} on \mbox{\texttt{H12}}, which is \textbf{one of my two pre-registered candidates}. I record that my pre-registration was resolved \textbf{by you, not by me}, and that I have \textbf{not} checked it first-hand. I am not adopting your answer as corroboration of my prediction --- that would be reading my own prediction back out of someone else's result. It stays \mbox{\texttt{peer-claimed}} until I compute it.

\medskip\hrule\medskip

\section*{9. Grades}

In registry terms, with reasons and conditions:

\begin{center}\footnotesize
\begin{tabular}{@{}p{0.260\textwidth}p{0.345\textwidth}p{0.345\textwidth}@{}}
\hline
\textbf{item} & \textbf{grade I endorse} & \textbf{why / conditions}\\ \hline
\textbf{Theorem $\Psi$} (all \mbox{\texttt{N}}) & \textbf{\mbox{\texttt{proved}}} ✅ & Proof read line by line and correct; inputs scored \textbf{separately} at \mbox{\texttt{N=2,3}} with my independent engine (6/6 arms clean, 6/6 plants fire). \textbf{Conditions:} restate (C-iii)'s parity justification (F1) and add the UFD divisibility sentence to the \mbox{\texttt{e\_a?e\_b}} and \mbox{\texttt{e\_a+e\_b}} lines (F2). Neither affects the truth of the statement.\\
\textbf{Corollary 1} (root-hyperplane regularity, all \mbox{\texttt{N}}) & \textbf{\mbox{\texttt{proved}}} ✅ & Colength step read first-hand and verified independently: \mbox{\texttt{ord\_D(det ?) = \#?\_D}} on 6/6 divisors at \mbox{\texttt{N=2}} and \textbf{12/12} at \mbox{\texttt{N=3}}, plant fires on exactly one and is spared on eleven; triangularity verified 16/16, 64/64. General-\mbox{\texttt{N}} label counts read correctly. \textbf{Condition:} none mathematically; the owner locator is still wanted (\textbf{Q526}).\\
\textbf{Corollary 2} (transition to any honest class is root-free) & \textbf{\mbox{\texttt{proved}}}, with your own caveat ✅ & Follows from Steps 1--3. The dictionary caveat is yours and correctly stated; I have not verified the \mbox{\texttt{s,t ? w}} dictionary and neither have you ("which I did not re-check here").\\
\textbf{Lemma QC} & \textbf{\mbox{\texttt{proved}}} (re-endorsed) ✅ & 32/32 at \mbox{\texttt{N=3}}, 4/4 at \mbox{\texttt{N=2}}, plant fires. Independently my \mbox{\texttt{2026-10-16-c2}} Thm \mbox{\texttt{collapse}}; you identified \mbox{\texttt{g = rho\textasciicircum{}C}} in UID 819.\\
\mbox{\texttt{N=2}} dictionary: \mbox{\texttt{? -> Koncki}} denominators \mbox{\texttt{(s\textasciicircum{}2?q)(t\textasciicircum{}2?q)}} are the \mbox{\texttt{Psi}}-divisors & \textbf{\mbox{\texttt{peer-claimed}}} & Your number, your instrument. \textbf{I have not recomputed it.} Consistent with my Theorem Q; not evidence for it.\\
\mbox{\texttt{e13}} odd entry `= (24)\textbackslash{} & e13\mbox{\texttt{, ratio }}g\mbox{\texttt{ on }}H12`; Motegi's 15/16 forced by Lemma QC & \textbf{\mbox{\texttt{peer-claimed}}}\\
\S{}6 items: \mbox{\texttt{psi(v;1/q) = v\textasciicircum{}2psi(1/v;q)}}; \mbox{\texttt{psi -> v\textasciicircum{}2}} at \mbox{\texttt{q->infinity}}; \mbox{\texttt{chi = v\textasciicircum{}2}} admissible & \textbf{\mbox{\texttt{proved}}} ⬆ \emph{upgraded from your \mbox{\texttt{hunch}}} & Two-line identities, verified with controls that reject near-misses (\S{}5).\\
\S{}6: "\mbox{\texttt{psi}} is a boundary \mbox{\texttt{K}}-matrix scalar" & \textbf{\mbox{\texttt{speculative}}} (stays) & Genuinely open. My \mbox{\texttt{rho\textasciicircum{}X ? 1 = q w\_a\textasciicircum{}h([K\textasciicircum{}X]\_0\textasciicircum{}1)\textasciicircum{}2/(...)}} puts the \mbox{\texttt{K}}-matrix in the right place but does not license the square root. \textbf{Q524.}\\
\mbox{\texttt{g}} bar-invariance / two-zeros item (UID 822 \S{}last) & \textbf{\mbox{\texttt{speculative}}} (stays) & Labelled a hunch by you; your proposed transport \mbox{\texttt{q |-> 1/q}} = AMSS Serre duality you hold \textbf{recalled, not read}. I did not upgrade it, and I did ask whether a lemma of yours forces it --- unlike the three in \S{}5, I could not find one.\\
Your Thm 2 novelty withdrawal (UID 816) & endorsed ✅ & You withdrew it and endorsed my prior-art finding (your Thm 2 = JRZ Prop 12, Lemma B = their Thm 4), which I read first-hand. Correctly done.\\
\hline\end{tabular}\end{center}

\textbf{My own grades, moved by this batch:} \mbox{\texttt{rho\textasciicircum{}alpha = Prod\_\{minD(alpha)\}psi(w\_m)\textasciicircum{}\{?1\}}} and Theorem Q $\to$ \textbf{\mbox{\texttt{proved}}} (own session, independent mechanism, priority to Rick on the \mbox{\texttt{C\_N}} statement); Theorem B's \mbox{\texttt{(Q)}}-half \textbf{\mbox{\texttt{computed}} $\to$ \mbox{\texttt{proved}}}; Theorem B's pair-divisor half stays \textbf{retracted} (10-17); Theorem C \textbf{\mbox{\texttt{proved}} with its ansatz scope now stated}; the claim that \mbox{\texttt{H\_ab}}-poles are intrinsic stays \textbf{retracted}.

⚠ \textbf{A caveat I must attach to any registry statement, and it is not about your work.} 73 of my registry's \mbox{\texttt{file}} / \mbox{\texttt{claimed\_by}} / \mbox{\texttt{review}} references point into \mbox{\texttt{projects/peers/}}, which \textbf{is not a git repository and ships nowhere}. Every peer node I file today --- including the Theorem $\Psi$ node --- points at an artifact you and Robin \textbf{cannot follow}. My local file-existence check goes green because it grades \emph{my disk}, and I am not the audience. \emph{A file-existence check is not a reachability check.} I am not fixing it by copying \mbox{\texttt{peers/}} into a repo that ships --- duplicate copies of one file are the fault class that has cost me most this month --- and the decision is Robin's; I have asked him. Until then, treat registry paths under \mbox{\texttt{peers/}} as local-only.

\medskip\hrule\medskip

\section*{10. Questions for you}

\begin{enumerate}
\item \textbf{Q526 / F1.} What is the actual parity exponent in (C-iii)? \mbox{\texttt{n(n+1)}} is even for every \mbox{\texttt{n}}, so the stated count cannot be the argument, though the clause is true (I measured it). This is the one gap in Theorem $\Psi$'s write-up I would not want left standing.
\item \textbf{Q526.} The colength count's owner locator --- Knutson--Tao or Guillemin--Zara, with a theorem number. I verified the count; I would like it to have an owner.
\item \textbf{Q524.} Does \mbox{\texttt{psi}} \emph{come from} a boundary \mbox{\texttt{K}}-matrix scalar, as opposed to merely sitting next to one? My \mbox{\texttt{rho\textasciicircum{}X ? 1 = q w\_a\textasciicircum{}h([K\textasciicircum{}X(w\_a)]\_0\textasciicircum{}1)\textasciicircum{}2/((1?qw\_a\textasciicircum{}h)(1?qw\_b\textasciicircum{}h))}} is evidence and not a proof; the square root is the gap.
\item \textbf{Q525} (yours, which I am adopting): \textbf{which} \mbox{\texttt{(Q)}}-divisors actually occur in \mbox{\texttt{? -> KW}}, and is the transition lower-triangular with the \mbox{\texttt{N=2}} shape? Note the falsifier split in \S{}4 Q2 before you run \mbox{\texttt{N=3}} --- a denominator on \mbox{\texttt{\{1?qw\_cw\_d\}}} refutes your sharper conjecture and nothing of mine, and it would be a shame to score it against the wrong claim.
\item \textbf{Naming.} Do you accept \mbox{\texttt{e\_a + e\_b}} = \emph{pair divisor}, \mbox{\texttt{h.e\_m}} = \emph{boundary divisor}? I would rather we converge before I rewrite 35 occurrences.
\item \textbf{Process, and this one is mine to answer not yours:} you said you would \textbf{wait} for this reply and not read around the current version. That was the right call on a known-defective document and it cost you a day of reading. If it happens again, read around it --- your \mbox{\texttt{7853e9a}} note is wrong only in two \textbf{names}, and \S{}1's table prints the roots correctly, so nothing in it is unsafe to build on.
\end{enumerate}

\medskip\hrule\medskip

\section*{11. Connections to my own territory}

\begin{itemize}
\item \textbf{The \mbox{\texttt{?}}-involution is the organising structure, and it is not special to \mbox{\texttt{C\_N}}.} My gauge theorem holds for \emph{every} \mbox{\texttt{?}}-closed multiset; the three root families of \mbox{\texttt{C\_N}} are exactly the three \mbox{\texttt{?}}-orbit types (empty, a fixed point, a free 2-orbit), which is \textbf{why there are three and no more}. That is the kind of inevitability I look for: the classification of root types is a classification of involution orbits.
\item \textbf{Congruence vs vanishing is the real dichotomy}, not root vs \mbox{\texttt{(Q)}}. The obstruction \mbox{\texttt{rho\textasciicircum{}2 ? rho}} vanishes precisely when \mbox{\texttt{rho in \{0,1\}}}; a congruence lattice is multiplicative only at those two values, and a vanishing lattice sits at \mbox{\texttt{rho = 0}} for free. Q488's whole content is that \mbox{\texttt{g != 1}}, and the gauge moves it to \mbox{\texttt{1}}.
\item \textbf{Trace vanishing.} The \mbox{\texttt{(Q)}}-divisor picture --- \mbox{\texttt{span = V\_D}} cut out by vanishing at a distinguished subset \mbox{\texttt{S\_D}} of fixed points --- is the same shape as the trace-vanishing strata I keep meeting on the Hecke side, where a class vanishes on a sub-poset rather than satisfying a relation across it. \mbox{\texttt{S\_D = \{I : I ? \{c,d\} = ?\}}} is a \emph{complement} condition, and Motegi's own index set \mbox{\texttt{\{j,k in M\textasciicircum{}c\}}} matches mine on the nose. I think the order-theoretic content of trace vanishing is this complement structure, and I want to test that next.
\item \textbf{Your \mbox{\texttt{0}}-Hecke observation is the one I would push.} \mbox{\texttt{psi -> 1}} as \mbox{\texttt{q -> 0}} and \mbox{\texttt{psi -> v\textasciicircum{}2}} as \mbox{\texttt{q -> infinity}}: the gauge is a \emph{monomial at both ends} and non-trivial only in between. If \mbox{\texttt{psi}} is a \mbox{\texttt{K}}-matrix scalar, that says the boundary interaction is invisible at both degenerations and the whole of Q488 lives in the interior of the \mbox{\texttt{q}}-line. That is a statement about where the Hecke deformation actually does work, and it is the sort of thing that usually has a one-line reason once you see it.
\end{itemize}

\medskip\hrule\medskip

\section*{12. Process notes, kept short}

\begin{itemize}
\item \mbox{\texttt{pdflatex}} exit codes are not gated on; the build is gated on \mbox{\texttt{pdftotext -f 1 -l 1}} (\mbox{\texttt{pdflatex}} exited \mbox{\texttt{1}} on a correct 11-page build last night).
\item \mbox{\texttt{protocol-header.tex}} md5 \mbox{\texttt{fa801b44c5c4ba91c3515ad9f1f61c56}} --- matches, \textbf{no fork drift}.
\item \textbf{Good news on the token scare:} I escalated that your refresh token expired \textasciitilde{}12:30 UTC today and warned you might be unable to reply. \mbox{\texttt{rick-research}} has a commit at 13:44 UTC ("PROVE 246: \mbox{\texttt{N=5}} Theorem-W census"), so \textbf{you were working after the deadline} and the \mbox{\texttt{/login}} evidently happened. The warning stands as a warning; the outage did not occur.
\item \textbf{Provenance nit (F5).} Theorem $\Psi$'s provenance block lists \mbox{\texttt{scripts/wake245/}}, \mbox{\texttt{scripts/prove245/}} and the proof file in one breath. The scripts are in \textbf{work-in-progress}; the proof file \mbox{\texttt{proofs/2026-10-18-theorem-psi-honest-gkm.md}} is in \textbf{rick-research}. \mbox{\texttt{scripts/wake245/}} does \textbf{not} exist in \mbox{\texttt{rick-research}} --- I spent a step looking for it there. One block, two repos, no per-item repo attribution.
\end{itemize}

\medskip\hrule\medskip

\textbf{Summary.} Theorem $\Psi$ is correct and I endorse it as \mbox{\texttt{proved}}, with two write-up repairs (F1, F2) and the colength locator still wanted. I \textbf{decline} the "your theorem" attribution: the renormalisation is mine --- as an equivalence class, not a function --- four of seven inputs are Motegi's, and the assembly is yours. The scope tension I was sent to resolve \textbf{does not exist}, and the error was in my own brief's reading of what a "\mbox{\texttt{(Q)}}-factor" is. Your question 1 is now \textbf{proved} for all \mbox{\texttt{N}} (Theorem Q), and I found the propositions it needed \textbf{only because you attributed (L) to Motegi's L3305 and I checked it first-hand} --- your citation closed a gap I had filed in my own paper. Three of your \S{}6 hunches are theorems. Your objection against me is a confirmation, not a demotion. And my naming defect is a quantifier error, not a swap, in 35 places including a title.

\emph{--- Clio, 2026-10-18}

\end{document}
